\section{Additional invariance and sensitivity results}
\subsection{Extension to regular connected curved interfaces}
\label{sec:curved}
Single-horizon rigidity uses analyticity and the gradient structure.
The reduction to the finite balances \eqref{eq:fiber} requires planes.
For a curved example, call a pair active if it shares a strict
two-class face. Suppose each active pair's full equality set
$H_{ab}=\{m_{ba}=0\}$ is a connected regular analytic hypersurface
and each boundary of a strict class cell is covered by the equality
sets of active pairs involving that class. Under these assumptions
preservation of every strict cell at one positive horizon is
equivalent to $F\cdot\nabla m_{ba}=0$ on each such equality set and to
preservation for all horizons.

Indeed, preservation maps a relatively open part of $H_{ab}$
into $H_{ab}$. The restriction of $m_{ba}\circ\Phi_T$ to that
connected analytic hypersurface is therefore zero. The same
compact-superlevel argument supplies an equilibrium there. Flatten
the hypersurface in a local analytic graph chart at the equilibrium
with orthonormal derivative. The linearization in those coordinates
is still symmetric, because the field vanishes at the equilibrium
and its derivative transforms by orthogonal conjugacy. The normal
Taylor calculation from Lemma~\ref{lem:sampled-tangency} now gives tangency
near that point, and analyticity extends $F\cdot\nabla m_{ba}=0$
along the whole hypersurface. Conversely, a first contact with the
boundary of the initial cell would lie on one of the invariant
hypersurfaces involving that class. A strict initial query does
not lie on that hypersurface, so backward uniqueness excludes
contact at finite time.

For example, if the prototypes of $b$ have strictly larger
projections onto some unit direction than all the prototypes of
$a$, then $H_{ab}$ is a connected graph over the orthogonal plane.
The directional derivative of $m_{ba}$ is the difference of the
two projected posterior means and is bounded below by the positive
projection gap. Hence there is one root on each parallel line,
and the implicit function theorem gives an analytic graph.
Interleaved classes in the core-only benchmark need not satisfy
the connectedness or regularity assumptions.

\subsection{Sensitivity to changes in the driving mixture}
\label{sec:sensitivity}
Fix the affine readout and let $F_\varepsilon$ be a smooth
one-parameter family of driving fields with $F_0$ preserving the
partition. Consider a compact patch $K$ in a strict two-class
facet with normal $n$. For $z\in K$, write $z_s=\Phi_s^0 z$ and
define the normal coefficients
\[
 \ell_z(s)=n^\top DF_0(z_s)n,\qquad
 b_z(s)=n^\top \partial_\varepsilon F_\varepsilon(z_s)\big|_0.
\]
At fixed time $T$, the initial points transported into the facet
have normal coordinate $r_\varepsilon(z)$ relative to $z$, where
\begin{equation}
 r_\varepsilon(z)=-\varepsilon I_T(z)+o(\varepsilon),
 \qquad
 I_T(z)=\int_0^T
       e^{-\int_0^s \ell_z(v)\,\dd v}\,b_z(s)\,\dd s .
 \label{eq:sensitivity}
\end{equation}
To derive this, the tangential derivatives of the base normal
velocity vanish identically on the plane. The first normal
variation of a perturbed trajectory initialized at $z$ therefore
obeys $\dot\eta=\ell_z\eta+b_z$ with $\eta(0)=0$.
The derivative of the terminal normal position with respect to
its initial normal coordinate is $\exp(\int_0^T\ell_z)$.
It is positive, so solving for the initial normal coordinate at
which the terminal coordinate vanishes gives \eqref{eq:sensitivity}
by the implicit function theorem. The formula holds uniformly on
$K$ by compactness. The unperturbed trajectory stays in the
relative interior of the facet because the base flow preserves
the closures of both adjacent cells and their intersection in
both time directions.

For a continuous feature density $\rho$, the probability of a
label switch with initial condition in a sufficiently thin
normal strip over $K$ is consequently
$|\varepsilon|\int_K\rho(z)|I_T(z)|\,\dd S+o(|\varepsilon|)$.
This is obtained by integrating across the interval between the
original and pulled-back interfaces. It describes a fixed finite
horizon under a fixed readout, whereas the earlier boundary flux
expands in time. Summing such formulas globally requires control
near junctions and at infinity.
